%=======================================================
% CHAPTER 2: HAMMING CODES
%=======================================================

\begin{frame}[plain]
  \begin{center}\bfseries\centering
    \includegraphics[scale=1]{assets/branding/opa0}
    \vspace{-5.75cm}\\
    {\adjustbox{scale=2.15}{\color{TealDarkCyan}\;III.\;Hamming Codes}}
    \bigskip\bigskip
    {\ }\vspace{4cm}\\
    \bigskip
  \end{center}
\end{frame}

%--- NEW: Historical Motivation ---
\begin{frame}{\textbf{III. Hamming Codes}}{Hamming (1950): The Original Problem}
\small
Richard Hamming worked at Bell Labs in 1947. The relay computer he used ran \textbf{punched-card jobs} overnight — and a single wrong hole would crash the job.

\medskip
\textbf{His frustration:} ``If the machine can detect an error, why can't it fix it?''

\medskip
\defbox[Hamming's insight (1950)]{
	Add \emph{redundant} parity bits so that:
	\begin{itemize}
		\item the syndrome identifies the \textbf{exact position} of the error, and
		\item \textbf{no syndrome is wasted} — each possible syndrome maps to exactly one error pattern.
	\end{itemize}
}

\medskip
This requires $n + 1 = 2^r$ syndromes for $n$ positions, giving $n = 2^r - 1$.
%\begin{columns}[T,onlytextwidth]
%  \begin{column}{.54\textwidth}
%    Richard Hamming worked at Bell Labs in 1947. The relay computer he used ran \textbf{punched-card jobs} overnight — and a single wrong hole would crash the job.
%
%    \medskip
%    \textbf{His frustration:} ``If the machine can detect an error, why can't it fix it?''
%
%    \medskip
%    \defbox[Hamming's insight (1950)]{
%    Add \emph{redundant} parity bits so that:
%    \begin{itemize}
%      \item the syndrome identifies the \textbf{exact position} of the error, and
%      \item \textbf{no syndrome is wasted} — each possible syndrome maps to exactly one error pattern.
%    \end{itemize}
%    }
%
%    \medskip
%    This requires $n + 1 = 2^r$ syndromes for $n$ positions, giving $n = 2^r - 1$.
%  \end{column}
%  \begin{column}{.43\textwidth}
%    \begin{tikzpicture}[scale=0.88]
%      % Show how syndromes map to errors
%      \node[font=\small\bfseries, TealDarkCyan] at (1.5,4.0)
%        {Syndrome $\to$ action};
%      \draw[TealDarkCyan, thick] (0,3.6)--(3.0,3.6);
%
%      % No error
%      \node[draw=TealDarkCyan!40, rounded corners=3pt, fill=TealDarkCyan!8,
%            minimum width=2.8cm, minimum height=0.38cm, font=\scriptsize]
%        at (1.5,3.25) {$000$ $\to$ no error};
%      % Errors at positions
%      \foreach \i/\bits in {1/001, 2/010, 3/011, 4/100, 5/101, 6/110, 7/111} {
%        \pgfmathsetmacro{\yy}{3.25 - \i*0.42}
%        \node[draw=AccentCyan!60, rounded corners=3pt, fill=AccentCyan!10,
%              minimum width=2.8cm, minimum height=0.38cm, font=\scriptsize]
%          at (1.5,\yy) {$\bits$ $\to$ flip bit $\i$};
%      }
%    \end{tikzpicture}
%  \end{column}
%\end{columns}
\end{frame}

\begin{frame}{\textbf{III. Hamming Codes}}{The design problem}
\small
\begin{block}{Goal}
Construct binary linear codes that correct exactly one error and use all possible syndromes.
\end{block}
\begin{itemize}
  \item A received word has either no error or one error in position $i$.
  \item Thus there are $1+n$ elementary cases.
  \item With $r$ parity checks, there are $2^r$ syndromes.
\end{itemize}
\[
  1+n=2^r
  \quad\Rightarrow\quad
  n=2^r-1.
\]
\end{frame}

\begin{frame}{\textbf{III. Hamming Codes}}{The parity-check matrix}
\small
\defbox[Binary Hamming matrix]{
For $r\ge2$, let $H_r$ be the $r\times(2^r-1)$ matrix whose columns are all nonzero vectors of $\F_2^r$.
}
\[
  H_3=
  \begin{bmatrix}
  1&0&1&0&1&0&1\\
  0&1&1&0&0&1&1\\
  0&0&0&1&1&1&1
  \end{bmatrix}.
\]
\begin{itemize}
  \item Column labels are nonzero binary syndromes.
  \item Every nonzero syndrome identifies a unique column.
\end{itemize}
\end{frame}

\begin{frame}{\textbf{III. Hamming Codes}}{Definition}
\small
\defbox[Binary Hamming code]{
The binary Hamming code of order $r$ is
\[
  \operatorname{Ham}(r,2)=\ker(H_r)\le\F_2^{2^r-1}.
\]
}
\begin{itemize}
  \item Length: $n=2^r-1$.
  \item Redundancy: $r$.
  \item Dimension: $k=n-r=2^r-1-r$.
\end{itemize}
\[
  \operatorname{Ham}(r,2)\text{ has parameters }[2^r-1,2^r-1-r,3]_2.
\]
\end{frame}

\begin{frame}{\textbf{III. Hamming Codes}}{Dimension}
\small
\thmbox[Dimension count]{
The matrix $H_r$ has rank $r$, hence
\[
  \dim\ker(H_r)=2^r-1-r.
\]
}
%\begin{proof}
%The standard basis vectors $e_1,\ldots,e_r$ occur among the nonzero columns of $H_r$, so the row rank is $r$.
%\end{proof}
%\begin{block}{Interpretation}
%The parity checks are independent; no check is redundant.
%\end{block}
\end{frame}

\begin{frame}{\textbf{III. Hamming Codes}}{Minimum distance}
\small
\thmbox[Distance]{
The binary Hamming code $\operatorname{Ham}(r,2)$ has minimum distance $3$.
}
%\begin{proof}
%No column of $H_r$ is zero, so $d\ge2$. No two columns are equal over $\F_2$, so no two are dependent, hence $d\ge3$. But $u+v+(u+v)=0$, so three suitable columns are dependent.
%\end{proof}
\begin{itemize}
  \item Therefore Hamming codes correct one error.
  \item They detect two errors.
\end{itemize}
\end{frame}

\begin{frame}{\textbf{III. Hamming Codes}}{Perfectness}
\small
\thmbox[Hamming codes are perfect]{
The binary Hamming code satisfies equality in the Hamming bound for $t=1$.
}
\[
  2^k(1+n)=2^{2^r-1-r}(1+2^r-1)=2^{2^r-1}=2^n.
\]
\begin{itemize}
  \item Every word of $\F_2^n$ is at distance at most $1$ from exactly one codeword.
  \item The syndrome is not merely a diagnostic: it is a complete decoder.
\end{itemize}
\end{frame}

%--- NEW: Perfect Tiling Visual ---
\begin{frame}{\textbf{III. Hamming Codes}}{Perfect Tiling: A Geometric View}
\small
\begin{columns}[T,onlytextwidth]
  \begin{column}{.48\textwidth}
    For the $[3,1,3]$ repetition code in $\F_2^3$: the two correction balls of radius 1 \textbf{tile the cube exactly} — no overlap, no gap.

    \medskip
    \begin{center}
    \begin{tikzpicture}[scale=0.80]
      % 3-cube edges
      \draw[gray!35,thick] (0,0)--(2,0); \draw[gray!35,thick] (0,0)--(1,0.9);
      \draw[gray!35,thick] (0,0)--(0.5,1.6); \draw[gray!35,thick] (2,0)--(3,0.9);
      \draw[gray!35,thick] (2,0)--(2.5,1.6); \draw[gray!35,thick] (1,0.9)--(3,0.9);
      \draw[gray!35,thick] (1,0.9)--(1.5,2.5); \draw[gray!35,thick] (0.5,1.6)--(2.5,1.6);
      \draw[gray!35,thick] (0.5,1.6)--(1.5,2.5); \draw[gray!35,thick] (3,0.9)--(3.5,2.5);
      \draw[gray!35,thick] (2.5,1.6)--(3.5,2.5); \draw[gray!35,thick] (1.5,2.5)--(3.5,2.5);
      % Ball around 000: {000,100,010,001} — TealDarkCyan fill
      \node[circle,draw=TealDarkCyan,fill=TealDarkCyan!20,minimum size=0.62cm,font=\tiny\bfseries]
        at (0,0)     {$000$};
      \node[circle,draw=TealDarkCyan,fill=TealDarkCyan!20,minimum size=0.62cm,font=\tiny\bfseries]
        at (2,0)     {$100$};
      \node[circle,draw=TealDarkCyan,fill=TealDarkCyan!20,minimum size=0.62cm,font=\tiny\bfseries]
        at (1,0.9)   {$010$};
      \node[circle,draw=TealDarkCyan,fill=TealDarkCyan!20,minimum size=0.62cm,font=\tiny\bfseries]
        at (0.5,1.6) {$001$};
      % Ball around 111: {111,011,101,110} — AccentCyan fill
      \node[circle,draw=AccentCyan,fill=AccentCyan!22,minimum size=0.62cm,font=\tiny\bfseries]
        at (3.5,2.5) {$111$};
      \node[circle,draw=AccentCyan,fill=AccentCyan!22,minimum size=0.62cm,font=\tiny\bfseries]
        at (1.5,2.5) {$011$};
      \node[circle,draw=AccentCyan,fill=AccentCyan!22,minimum size=0.62cm,font=\tiny\bfseries]
        at (2.5,1.6) {$101$};
      \node[circle,draw=AccentCyan,fill=AccentCyan!22,minimum size=0.62cm,font=\tiny\bfseries]
        at (3,0.9)   {$110$};
      % Codeword labels
      \node[below=3pt,font=\tiny,TealDarkCyan] at (0,0)     {\textbf{c=000}};
      \node[above=3pt,font=\tiny,AccentCyan]   at (3.5,2.5) {\textbf{c=111}};
    \end{tikzpicture}
	\end{center}
  \end{column}
  \begin{column}{.49\textwidth}
    \textcolor{TealDarkCyan}{\textbf{Why "perfect"?}}
    \begin{itemize}
      \item Each ball has 4 elements: 1 codeword + 3 neighbors
      \item Two balls cover all $2^3 = 8$ points of $\F_2^3$
      \item $2 \times 4 = 8$ — the Hamming bound holds with \emph{equality}
    \end{itemize}

%    \medskip
%    \defbox[For $\operatorname{Ham}(r,2)$]{
%    $2^{2^r-1-r} \times (1 + 2^r - 1) = 2^{2^r-1}$ — all of $\F_2^{2^r-1}$ is tiled.\\[2pt]
%    Every received word is within distance 1 of \textbf{exactly one} codeword.
%    }

    \medskip
    Hamming codes are the \emph{only} linear perfect 1-error-correcting codes (apart from trivial cases).
  \end{column}
\end{columns}
\end{frame}

\begin{frame}{\textbf{III. Hamming Codes}}{Syndrome decoding}
\small
\defbox[Decoder]{
Given $v\in\F_2^n$, compute $s=H_rv^T$.
}
\begin{itemize}
  \item If $s=0$, accept $v$.
  \item If $s\ne0$, find the unique column $H_i=s$.
  \item Correct by flipping coordinate $i$.
\end{itemize}
%\begin{block}{Reason}
%For a single error $e_i$, the syndrome is $H_re_i^T=H_i$.
%\end{block}
\end{frame}

%--- NEW: Full Syndrome Decoder Table ---
\begin{frame}{\textbf{III. Hamming Codes}}{The Complete Syndrome Table for $[7,4,3]$}
\small
With $H_3$ as written, the syndrome $(s_1,s_2,s_3)^T$ equals column $i$ when bit $i$ is flipped.

\medskip
\begin{center}
\begin{tikzpicture}[scale=0.82]
  % Header
  \node[font=\scriptsize\bfseries,TealDarkCyan] at (0.0,3.5) {Syndrome %$(s_1\,s_2\,s_3)$
  };
  \node[font=\scriptsize\bfseries,TealDarkCyan] at (2.5,3.5) {Error pattern $e$};
  \node[font=\scriptsize\bfseries,TealDarkCyan] at (5.2,3.5) {Action};
  \draw[TealDarkCyan!50,thick] (-0.8,3.25)--(6.0,3.25);
  % Row 0: no error
  \node[draw=TealDarkCyan!30,fill=TealDarkCyan!8,rounded corners=2pt,
        minimum width=2.0cm,minimum height=0.34cm,font=\scriptsize]
    at (0.0,3.0) {$(0,0,0)$};
  \node[font=\scriptsize] at (2.5,3.0) {$0000000$};
  \node[font=\scriptsize,TealDarkCyan] at (5.2,3.0) {no error — accept};
  % Rows 1-7: each error position
  \node[draw=AccentCyan,fill=AccentCyan!12,rounded corners=2pt,
        minimum width=2.0cm,minimum height=0.34cm,font=\scriptsize]
    at (0.0,2.6) {$(1,0,0)$};
  \node[font=\scriptsize] at (2.5,2.6) {$1000000$};
  \node[font=\scriptsize,AccentCyan] at (5.2,2.6) {flip bit 1};

  \node[draw=AccentCyan,fill=AccentCyan!12,rounded corners=2pt,
        minimum width=2.0cm,minimum height=0.34cm,font=\scriptsize]
    at (0.0,2.2) {$(0,1,0)$};
  \node[font=\scriptsize] at (2.5,2.2) {$0100000$};
  \node[font=\scriptsize,AccentCyan] at (5.2,2.2) {flip bit 2};

  \node[draw=AccentCyan,fill=AccentCyan!12,rounded corners=2pt,
        minimum width=2.0cm,minimum height=0.34cm,font=\scriptsize]
    at (0.0,1.8) {$(1,1,0)$};
  \node[font=\scriptsize] at (2.5,1.8) {$0010000$};
  \node[font=\scriptsize,AccentCyan] at (5.2,1.8) {flip bit 3};

  \node[draw=AccentCyan,fill=AccentCyan!12,rounded corners=2pt,
        minimum width=2.0cm,minimum height=0.34cm,font=\scriptsize]
    at (0.0,1.4) {$(0,0,1)$};
  \node[font=\scriptsize] at (2.5,1.4) {$0001000$};
  \node[font=\scriptsize,AccentCyan] at (5.2,1.4) {flip bit 4};

  \node[draw=AccentCyan,fill=AccentCyan!12,rounded corners=2pt,
        minimum width=2.0cm,minimum height=0.34cm,font=\scriptsize]
    at (0.0,1.0) {$(1,0,1)$};
  \node[font=\scriptsize] at (2.5,1.0) {$0000100$};
  \node[font=\scriptsize,AccentCyan] at (5.2,1.0) {flip bit 5};

  \node[draw=AccentCyan,fill=AccentCyan!12,rounded corners=2pt,
        minimum width=2.0cm,minimum height=0.34cm,font=\scriptsize]
    at (0.0,0.6) {$(0,1,1)$};
  \node[font=\scriptsize] at (2.5,0.6) {$0000010$};
  \node[font=\scriptsize,AccentCyan] at (5.2,0.6) {flip bit 6};

  \node[draw=AccentCyan,fill=AccentCyan!12,rounded corners=2pt,
        minimum width=2.0cm,minimum height=0.34cm,font=\scriptsize]
    at (0.0,0.2) {$(1,1,1)$};
  \node[font=\scriptsize] at (2.5,0.2) {$0000001$};
  \node[font=\scriptsize,AccentCyan] at (5.2,0.2) {flip bit 7};
\end{tikzpicture}
\end{center}

\vspace{-2pt}
\begin{itemize}
  \item All $2^3 = 8$ syndromes are used — confirming \textbf{perfectness}
  \item Decoding is a single table look-up — $O(1)$ hardware
\end{itemize}
\end{frame}

\begin{frame}{\textbf{III. Hamming Codes}}{Binary index convention}
\small
\defbox[Indexed columns]{
If the $i$-th column of $H_r$ is the binary expansion of $i$, then the nonzero syndrome is literally the error position.
}
\[
  s=(1,0,1)^T\quad\leadsto\quad i=5
\]
depending on the chosen bit order.
\begin{itemize}
  \item This is a notational convention, not an invariant.
  \item It explains the classical parity-bit positions $1,2,4,\ldots$.
\end{itemize}
\end{frame}

\begin{frame}{\textbf{III. Hamming Codes}}{The classical [7,4,3] code}
\small
\[
H=
\begin{bmatrix}
1&0&1&0&1&0&1\\
0&1&1&0&0&1&1\\
0&0&0&1&1&1&1
\end{bmatrix}.
\]
\begin{itemize}
  \item There are $7$ nonzero syndromes.
  \item Each syndrome names one error coordinate.
  \item The kernel has dimension $4$.
\end{itemize}
\[
  \operatorname{Ham}(3,2)=[7,4,3]_2.
\]
\end{frame}

%--- NEW: Parity Bit Layout ---
\begin{frame}{\textbf{III. Hamming Codes}}{The $[7,4,3]$ Code: Parity Bit Positions}
\small
Bits are numbered $1$ through $7$. Positions that are powers of 2 ($1, 2, 4$) are \textcolor{AccentCyan}{\textbf{parity bits}}; the rest carry \textcolor{TealDarkCyan}{\textbf{data}}.

\medskip
\begin{center}
\begin{tikzpicture}[scale=0.88]
  % Seven position boxes
  \foreach \i/\type in {1/par, 2/par, 3/data, 4/par, 5/data, 6/data, 7/data} {
    \pgfmathsetmacro{\xp}{(\i-1)*1.05}
    \ifnum\i=1
      \node[draw=AccentCyan, thick, fill=AccentCyan!20, rounded corners=4pt,
            minimum width=0.9cm, minimum height=0.9cm, font=\small\bfseries]
        (b\i) at (\xp,0) {$p_1$};
    \else\ifnum\i=2
      \node[draw=AccentCyan, thick, fill=AccentCyan!20, rounded corners=4pt,
            minimum width=0.9cm, minimum height=0.9cm, font=\small\bfseries]
        (b\i) at (\xp,0) {$p_2$};
    \else\ifnum\i=4
      \node[draw=AccentCyan, thick, fill=AccentCyan!20, rounded corners=4pt,
            minimum width=0.9cm, minimum height=0.9cm, font=\small\bfseries]
        (b\i) at (\xp,0) {$p_4$};
    \else
      \node[draw=TealDarkCyan, thick, fill=TealDarkCyan!12, rounded corners=4pt,
            minimum width=0.9cm, minimum height=0.9cm, font=\small\bfseries]
        (b\i) at (\xp,0) {$d_{\i}$};
    \fi\fi\fi
    \node[font=\tiny, TealDarkCyan!70] at (\xp,-0.62) {bit \i};
  }
\end{tikzpicture}
\end{center}

\medskip
\textbf{Parity check coverage:}
\begin{center}
\begin{tabular}{c|l|l}
  \textbf{Parity bit} & \textbf{Checks bits} & \textbf{Syndrome row} \\\hline
  $p_1$ (bit 1) & 1, 3, 5, 7 & $H$ row 1: $(1,0,1,0,1,0,1)$ \\
  $p_2$ (bit 2) & 2, 3, 6, 7 & $H$ row 2: $(0,1,1,0,0,1,1)$ \\
  $p_4$ (bit 4) & 4, 5, 6, 7 & $H$ row 3: $(0,0,0,1,1,1,1)$ \\
\end{tabular}
\end{center}

\medskip
\begin{itemize}
  \item Each parity bit checks all positions whose index has a $1$ in the corresponding binary place
  \item The 3-bit syndrome tells us the \emph{binary address} of the flipped bit
\end{itemize}
\end{frame}

\begin{frame}{\textbf{III. Hamming Codes}}{A systematic generator matrix}
\small
One systematic form for the $[7,4,3]$ Hamming code is
\[
G=
\begin{bmatrix}
1&0&0&0&0&1&1\\
0&1&0&0&1&0&1\\
0&0&1&0&1&1&0\\
0&0&0&1&1&1&1
\end{bmatrix}.
\]
\begin{itemize}
  \item The first four coordinates are message symbols.
  \item The last three are parity symbols.
  \item This $G$ is equivalent to the kernel of a suitable $H$.
\end{itemize}
\end{frame}

\begin{frame}{\textbf{III. Hamming Codes}}{Encoding example}
\small
For $m=(1,0,1,1)$ and the systematic matrix $G$,
\[
  c=mG=(1,0,1,1,0,0,1).
\]
\begin{itemize}
  \item The parity bits are linear forms in the message bits.
  \item A single flipped bit changes the syndrome to the corresponding column.
  \item The correction rule does not require knowing the transmitted message.
\end{itemize}
\end{frame}

%--- NEW: Encoding and Error Correction Flow ---
\begin{frame}{\textbf{III. Hamming Codes}}{Encoding \& Error Correction}
\small
\begin{center}
\begin{tikzpicture}[
  box/.style={draw=TealDarkCyan,thick,rounded corners=4pt,
              minimum width=2.0cm,minimum height=0.65cm,
              fill=TealDarkCyan!10,font=\scriptsize,align=center},
  arr/.style={-Stealth,thick,TealDarkCyan},
  earr/.style={-Stealth,thick,AccentCyan,dashed}
]
  % Nodes
  \node[box] (msg)  at (0,0)   {Message\\$m=(1,0,1,1)$};
  \node[box] (enc)  at (3,0)   {Encoder\\$c=mG$};
  \node[box] (cw)   at (6,0)   {Codeword\\$c=(1,0,1,1,0,1,0)$};
  \node[box,fill=AccentCyan!12,draw=AccentCyan] (ch) at (9,0)
    {Channel\\flip bit 5};
  \node[box,fill=AccentCyan!8,draw=AccentCyan]  (rv)  at (9,-1.5)
    {Received\\$v=(1,0,1,1,\mathbf{1},1,0)$};
  \node[box] (syn)  at (6,-1.5) {Syndrome\\$s=Hv^T=(1,0,1)^T$};
  \node[box] (lkp)  at (3,-1.5) {Look up:\\col 5 of $H$};
  \node[box] (fix)  at (0,-1.5) {Correct:\\$\hat{c}=c$\;\checkmark};
  % Arrows
  \draw[arr]  (msg)--(enc);
  \draw[arr]  (enc)--(cw);
  \draw[arr]  (cw)--(ch);
  \draw[earr] (ch)--(rv);
  \draw[arr]  (rv)--(syn);
  \draw[arr]  (syn)--(lkp);
  \draw[arr]  (lkp)--(fix);
  % Labels
  \node[above=2pt,font=\tiny,TealDarkCyan] at (1.5,0)   {$mG$};
  \node[above=2pt,font=\tiny,AccentCyan]   at (9,-0.75) {noise};
  \node[above=2pt,font=\tiny,TealDarkCyan] at (7.5,-1.5) {$Hv^T$};
\end{tikzpicture}
\end{center}

\begin{columns}[T,onlytextwidth]
  \begin{column}{.50\textwidth}
    \textbf{Key numbers:}
    \begin{itemize}
      \item Syndrome $s=(1,0,1)^T$ matches column 5 of $H_3$
      \item Flip bit 5: $(1,0,1,1,\mathbf{1},1,0)\to(1,0,1,1,0,1,0)$
      \item Recovered $\hat{c}$ is the original codeword (v)
    \end{itemize}
  \end{column}
  \begin{column}{.47\textwidth}
%    \begin{block}{The decoder never needs $m$}
	The decoder never needs $m$\\
    The syndrome alone identifies the error. Once we flip the right bit, the first 4 bits of the corrected $\hat{c}$ directly give $\hat{m}=(1,0,1,1)$.
%    \end{block}
  \end{column}
\end{columns}
\end{frame}

\begin{frame}{\textbf{III. Hamming Codes}}{Projective geometry}
\small
\defbox[Projective interpretation]{
The nonzero columns of $H_r$, modulo scalar multiplication, are the points of projective space
\[
  \mathbb{P}^{r-1}(\F_q).
\]
}
\begin{itemize}
  \item Over $\F_2$, each nonzero vector is its own projective point.
  \item Over $\F_q$, scalar multiples define the same projective point.
  \item Hamming codes are parity checks indexed by projective points.
\end{itemize}
\end{frame}

\begin{frame}{\textbf{III. Hamming Codes}}{The Fano plane}
\small
For $r=3$, the projective space $\mathbb{P}^2(\F_2)$ has $7$ points and $7$ lines.
\begin{itemize}
  \item The $7$ columns of $H_3$ are the points.
  \item A dependent triple of columns is a line of the Fano plane.
  \item Weight-$3$ codewords are exactly these lines.
\end{itemize}
%\begin{block}{Geometry}
The minimum distance $3$ is the statement that the smallest projective line has three points.
%\end{block}
\end{frame}

%--- NEW: Fano Plane Drawing ---
\begin{frame}{\textbf{III. Hamming Codes}}{The Fano Plane — Drawn}
\small
\begin{columns}[T,onlytextwidth]
  \begin{column}{.44\textwidth}
    \begin{center}
    \begin{tikzpicture}[scale=0.88]
      % Outer triangle edges (3 lines of the Fano plane)
      % Line {1,2,3}: P1(0,2) -- P3(-0.865,0.5) -- P2(-1.73,-1)
      \draw[TealDarkCyan, thick] (0,2) -- (-1.73,-1);
      % Line {1,4,5}: P1(0,2) -- P5(0.865,0.5) -- P4(1.73,-1)
      \draw[TealDarkCyan, thick] (0,2) -- (1.73,-1);
      % Line {2,4,6}: P2(-1.73,-1) -- P6(0,-1) -- P4(1.73,-1)
      \draw[TealDarkCyan, thick] (-1.73,-1) -- (1.73,-1);
      % Three "altitude" lines (through center P7=(0,0))
      % Line {1,7,6}: P1(0,2) -- P7(0,0) -- P6(0,-1)
      \draw[AccentCyan, thick] (0,2) -- (0,-1);
      % Line {2,7,5}: P2(-1.73,-1) -- P7(0,0) -- P5(0.865,0.5)
      \draw[AccentCyan, thick] (-1.73,-1) -- (0.865,0.5);
      % Line {3,7,4}: P3(-0.865,0.5) -- P7(0,0) -- P4(1.73,-1)
      \draw[AccentCyan, thick] (-0.865,0.5) -- (1.73,-1);
      % Inner circle through {3,5,6} (circumradius ≈ 1 centered at P7=(0,0))
      \draw[TealDarkCyan!55, thick] (0,0) circle (1.0);
      % Points — outer triangle (columns 1,2,4)
      \fill[TealDarkCyan] (0,2)      circle (3.5pt);
      \node[above=3pt, font=\scriptsize, TealDarkCyan]
        at (0,2)      {$1\ (100)$};
      \fill[TealDarkCyan] (-1.73,-1) circle (3.5pt);
      \node[left=3pt,  font=\scriptsize, TealDarkCyan]
        at (-1.73,-1) {$2\ (010)$};
      \fill[TealDarkCyan] (1.73,-1)  circle (3.5pt);
      \node[right=3pt, font=\scriptsize, TealDarkCyan]
        at (1.73,-1)  {$4\ (001)$};
      % Edge midpoints (columns 3,5,6)
      \fill[AccentCyan] (-0.865,0.5) circle (3.5pt);
      \node[left=3pt,  font=\scriptsize, AccentCyan]
        at (-0.865,0.5) {$3\ (110)$};
      \fill[AccentCyan] (0.865,0.5)  circle (3.5pt);
      \node[right=3pt, font=\scriptsize, AccentCyan]
        at (0.865,0.5)  {$5\ (101)$};
      \fill[AccentCyan] (0,-1)        circle (3.5pt);
      \node[below=3pt, font=\scriptsize, AccentCyan]
        at (0,-1)       {$6\ (011)$};
      % Center (column 7)
      \fill[TealDarkCyan!70] (0,0) circle (3.5pt);
      \node[right=4pt, font=\scriptsize, TealDarkCyan]
        at (0,0) {$7\ (111)$};
    \end{tikzpicture}
    \end{center}
  \end{column}
  \begin{column}{.53\textwidth}
    \textcolor{TealDarkCyan}{\textbf{The 7 lines (each has 3 points):}}
    \smallskip

    \begin{tabular}{c|l|l}
      \# & \textbf{Points} & \textbf{Columns sum to 0} \\\hline
      \rowcolor{TealDarkCyan!8}
      1 & $\{1,2,3\}$ & $100{+}010{=}110$ \\
      2 & $\{1,4,5\}$ & $100{+}001{=}101$ \\
      \rowcolor{TealDarkCyan!8}
      3 & $\{2,4,6\}$ & $010{+}001{=}011$ \\
      4 & $\{1,6,7\}$ & $100{+}011{=}111$ \\
      \rowcolor{TealDarkCyan!8}
      5 & $\{2,5,7\}$ & $010{+}101{=}111$ \\
      6 & $\{3,4,7\}$ & $110{+}001{=}111$ \\
      \rowcolor{TealDarkCyan!8}
      7 & $\{3,5,6\}$ & $110{+}101{=}011$ \\
    \end{tabular}

%    \bigskip
%    \begin{block}{Why minimum distance $= 3$}
%    Each line is a weight-$3$ codeword. No point lies alone (weight-$1$), no pair sums to zero (weight-$2$). Hence $d = 3$.
%    \end{block}
  \end{column}
\end{columns}
\end{frame}

\begin{frame}{\textbf{III. Hamming Codes}}{Lines give codewords}
\small
In $\mathbb{P}^2(\F_2)$, if $a,b,c$ are the three points on a line, then
\[
  a+b+c=0.
\]
\begin{itemize}
  \item The incidence vector of that line is a Hamming codeword.
  \item The $7$ lines give the $7$ words of weight $3$.
  \item Their complements give $7$ words of weight $4$.
\end{itemize}
\[
  W_{\operatorname{Ham}(3,2)}(Y)=1+7Y^3+7Y^4+Y^7.
\]
\end{frame}

\begin{frame}{\textbf{III. Hamming Codes}}{The simplex dual}
\small
\defbox[Simplex code]{
The row space of $H_r$ is the binary simplex code:
\[
  \operatorname{Sim}(r,2)=\operatorname{rowspan}(H_r).
\]
}
\begin{itemize}
  \item It has parameters $[2^r-1,r,2^{r-1}]_2$.
  \item Every nonzero word has weight $2^{r-1}$.
  \item It is dual to the Hamming code.
\end{itemize}
\[
  \operatorname{Ham}(r,2)^\perp=\operatorname{Sim}(r,2).
\]
\end{frame}

\begin{frame}{\textbf{III. Hamming Codes}}{MacWilliams check}
\small
For $r=3$, the simplex code has weight enumerator
\[
  W_{\operatorname{Sim}}(X,Y)=X^7+7X^3Y^4.
\]
Applying MacWilliams gives
\[
  W_{\operatorname{Ham}}(X,Y)=X^7+7X^4Y^3+7X^3Y^4+Y^7.
\]
\begin{itemize}
  \item The algebraic transform recovers the geometric weight distribution.
  \item Duality is visible both in matrices and in enumerators.
\end{itemize}
\end{frame}

\begin{frame}{\textbf{III. Hamming Codes}}{Q-ary Hamming codes}
\small
\defbox[Q-ary construction]{
Choose one nonzero representative from each one-dimensional subspace of $\F_q^r$ and use these vectors as columns of $H$.
}
\[
  n=\frac{q^r-1}{q-1}.
\]
\begin{itemize}
  \item The resulting code has parameters
  \[
    \left[\frac{q^r-1}{q-1},\frac{q^r-1}{q-1}-r,3\right]_q.
  \]
  \item It is again perfect for one-error correction.
\end{itemize}
\end{frame}

\begin{frame}{\textbf{III. Hamming Codes}}{Why projective representatives}
\small
Over $\F_q$, two proportional nonzero columns are linearly dependent.
\begin{itemize}
  \item To have $d\ge3$, no two columns may be proportional.
  \item Therefore columns must be projective points, not raw nonzero vectors.
  \item Taking all projective points uses the maximum possible number of columns.
\end{itemize}
%\begin{block}{Conclusion}
Hamming codes are the maximal parity-check configurations with no zero column and no repeated projective point.
%\end{block}
\end{frame}

\begin{frame}{\textbf{III. Hamming Codes}}{Q-ary perfectness}
\small
For the $q$-ary Hamming code,
\[
  n=\frac{q^r-1}{q-1},\qquad k=n-r.
\]
The radius-one sphere has size
\[
  1+n(q-1)=1+q^r-1=q^r.
\]
Hence
\[
  q^k(1+n(q-1))=q^{n-r}q^r=q^n.
\]
%\begin{block}{Perfect again}
Every received word lies at distance at most one from a unique codeword.
%\end{block}
\end{frame}

%%--- NEW: Hamming Family Parameters ---
%\begin{frame}{\textbf{III. Hamming Codes}}{The Hamming Family: Parameters at a Glance}
%\small
%\begin{center}
%\begin{tikzpicture}[scale=0.88]
%  % Table header
%  \node[font=\scriptsize\bfseries,TealDarkCyan] at (0.0,4.5) {$r$};
%  \node[font=\scriptsize\bfseries,TealDarkCyan] at (1.2,4.5) {$n=2^r{-}1$};
%  \node[font=\scriptsize\bfseries,TealDarkCyan] at (2.6,4.5) {$k=n{-}r$};
%  \node[font=\scriptsize\bfseries,TealDarkCyan] at (3.9,4.5) {Rate $k/n$};
%  \node[font=\scriptsize\bfseries,TealDarkCyan] at (5.3,4.5) {$t$ corrected};
%  \node[font=\scriptsize\bfseries,TealDarkCyan] at (6.7,4.5) {Use case};
%  \draw[TealDarkCyan!40,thick] (-0.4,4.25)--(7.5,4.25);
%  % r=2
%  \node[font=\scriptsize] at (0.0,3.9) {2};
%  \node[font=\scriptsize] at (1.2,3.9) {3};
%  \node[font=\scriptsize] at (2.6,3.9) {1};
%  \node[font=\scriptsize] at (3.9,3.9) {$33\%$};
%  \node[font=\scriptsize] at (5.3,3.9) {1};
%  \node[font=\scriptsize,gray] at (6.7,3.9) {tiny};
%  % r=3
%  \node[draw=TealDarkCyan!30,fill=TealDarkCyan!8,rounded corners=2pt,
%        minimum width=6.6cm,minimum height=0.34cm] at (3.35,3.45) {};
%  \node[font=\scriptsize\bfseries,TealDarkCyan] at (0.0,3.45) {3};
%  \node[font=\scriptsize\bfseries,TealDarkCyan] at (1.2,3.45) {7};
%  \node[font=\scriptsize\bfseries,TealDarkCyan] at (2.6,3.45) {4};
%  \node[font=\scriptsize\bfseries,TealDarkCyan] at (3.9,3.45) {$57\%$};
%  \node[font=\scriptsize\bfseries,TealDarkCyan] at (5.3,3.45) {1};
%  \node[font=\scriptsize\bfseries,TealDarkCyan] at (6.7,3.45) {[7,4,3]};
%  % r=4
%  \node[font=\scriptsize] at (0.0,3.0) {4};
%  \node[font=\scriptsize] at (1.2,3.0) {15};
%  \node[font=\scriptsize] at (2.6,3.0) {11};
%  \node[font=\scriptsize] at (3.9,3.0) {$73\%$};
%  \node[font=\scriptsize] at (5.3,3.0) {1};
%  \node[font=\scriptsize,gray] at (6.7,3.0) {medium};
%  % r=5
%  \node[draw=TealDarkCyan!30,fill=TealDarkCyan!8,rounded corners=2pt,
%        minimum width=6.6cm,minimum height=0.34cm] at (3.35,2.55) {};
%  \node[font=\scriptsize] at (0.0,2.55) {5};
%  \node[font=\scriptsize] at (1.2,2.55) {31};
%  \node[font=\scriptsize] at (2.6,2.55) {26};
%  \node[font=\scriptsize] at (3.9,2.55) {$84\%$};
%  \node[font=\scriptsize] at (5.3,2.55) {1};
%  \node[font=\scriptsize,gray] at (6.7,2.55) {medium};
%  % r=6
%  \node[font=\scriptsize] at (0.0,2.1) {6};
%  \node[font=\scriptsize] at (1.2,2.1) {63};
%  \node[font=\scriptsize] at (2.6,2.1) {57};
%  \node[font=\scriptsize] at (3.9,2.1) {$90\%$};
%  \node[font=\scriptsize] at (5.3,2.1) {1};
%  \node[font=\scriptsize,gray] at (6.7,2.1) {large};
%  % r=8 (ECC RAM extended)
%  \node[draw=AccentCyan!20,fill=AccentCyan!8,rounded corners=2pt,
%        minimum width=6.6cm,minimum height=0.34cm] at (3.35,1.65) {};
%  \node[font=\scriptsize,AccentCyan] at (0.0,1.65) {$*$};
%  \node[font=\scriptsize\bfseries,AccentCyan] at (1.2,1.65) {72};
%  \node[font=\scriptsize\bfseries,AccentCyan] at (2.6,1.65) {64};
%  \node[font=\scriptsize\bfseries,AccentCyan] at (3.9,1.65) {$89\%$};
%  \node[font=\scriptsize\bfseries,AccentCyan] at (5.3,1.65) {1 (SECDED)};
%  \node[font=\scriptsize\bfseries,AccentCyan] at (6.7,1.65) {ECC RAM};
%\end{tikzpicture}
%\end{center}
%
%\vspace{-4pt}
%{\footnotesize $*$ Extended Hamming code $[72,64,4]_2$: shortened and extended for byte-aligned memory.}
%
%\begin{itemize}
%  \item Rate $\to 1$ as $r \to \infty$ — Hamming codes become increasingly efficient
%  \item All correct exactly \textbf{one} bit error (single-error-correcting)
%\end{itemize}
%\end{frame}

\begin{frame}{\textbf{III. Hamming Codes}}{Extended Hamming codes}
\small
\defbox[Extension]{
Add one overall parity coordinate to force even total weight.
}
\begin{itemize}
  \item The binary $[7,4,3]$ code becomes an $[8,4,4]$ code.
  \item The new minimum distance is $4$.
  \item It detects three errors and corrects one error.
\end{itemize}
\[
  \widehat{\operatorname{Ham}}(3,2)=[8,4,4]_2.
\]
\end{frame}

\begin{frame}{\textbf{III. Hamming Codes}}{The [8,4,4] code}
\small
The extended Hamming code has weight enumerator
\[
  W(X,Y)=X^8+14X^4Y^4+Y^8.
\]
\begin{itemize}
  \item All codeword weights are divisible by $4$.
  \item The code is self-dual.
  \item It is the smallest doubly-even self-dual binary code.
\end{itemize}
%\begin{block}{Connection}
This code is the length-$8$ seed behind the $E_8$ lattice construction.
%\end{block}
\end{frame}

\begin{frame}{\textbf{III. Hamming Codes}}{Automorphisms}
\small
\defbox[Symmetry group]{
The automorphism group of the binary Hamming code is induced by invertible linear transformations of $\F_2^r$:
\[
  \Aut(\operatorname{Ham}(r,2))\cong \operatorname{GL}(r,2).
\]
}
\begin{itemize}
  \item A matrix in $\operatorname{GL}(r,2)$ permutes the nonzero columns of $H_r$.
  \item Hence it permutes coordinates and preserves the kernel.
  \item The code remembers the projective geometry.
\end{itemize}
\end{frame}

\begin{frame}{\textbf{III. Hamming Codes}}{Tanner graph}
\small
\begin{itemize}
  \item The Tanner graph of $H_r$ has $n=2^r-1$ variable nodes and $r$ check nodes.
  \item It is dense compared with modern LDPC codes.
  \item But it has an exceptionally simple exact decoder.
\end{itemize}
\begin{block}{Interpretation}
For Hamming codes, the syndrome is not iteratively refined; it is read directly as a projective point.
\end{block}
\end{frame}

\begin{frame}{\textbf{III. Hamming Codes}}{Failure modes}
\small
\begin{itemize}
  \item A two-error pattern has syndrome equal to the sum of two columns.
  \item In a binary Hamming matrix, that sum is another nonzero column.
  \item Therefore a two-error pattern may be miscorrected as a one-error pattern.
\end{itemize}
%\begin{block}{Moral}
Minimum distance $3$ guarantees single-error correction, not reliable detection after attempting correction.
%\end{block}
\end{frame}

%\begin{frame}{\textbf{III. Hamming Codes}}{SECDED}
%\small
%\defbox[SECDED]{
%Single-error correction, double-error detection is obtained by using the extended Hamming code.
%}
%\begin{itemize}
%  \item The extension raises the distance from $3$ to $4$.
%  \item One error is corrected by the Hamming syndrome.
%  \item Two errors are detected because parity and syndrome disagree in a recognizable way.
%\end{itemize}
%\end{frame}
%
%%--- NEW: ECC RAM Application ---
%\begin{frame}{\textbf{III. Hamming Codes}}{ECC Memory: Hamming in Your Computer}
%\small
%\begin{columns}[T,onlytextwidth]
%  \begin{column}{.54\textwidth}
%    \textcolor{TealDarkCyan}{\textbf{Error-correcting DRAM (ECC RAM)}}
%
%    \medskip
%    Standard server RAM uses the \textbf{extended Hamming code $[72, 64, 4]_2$}:
%    \begin{itemize}
%      \item $64$ data bits + $8$ check bits per 64-bit memory word
%      \item Corrects \emph{any single-bit} error automatically
%      \item Detects (but does not correct) any \emph{double-bit} error
%      \item Rate $= 64/72 \approx 89\%$ — very efficient
%    \end{itemize}
%
%    \medskip
%    \textcolor{TealDarkCyan}{\textbf{Why does it matter?}}
%    \begin{itemize}
%      \item Cosmic rays and thermal noise can flip a bit in DRAM roughly once per gigabyte per day
%      \item In a server with $1\,\text{TB}$ of RAM, that is $\sim 1000$ bit flips per day
%      \item ECC corrects each flip silently — the OS never sees it
%    \end{itemize}
%  \end{column}
%  \begin{column}{.43\textwidth}
%    \begin{tikzpicture}[scale=0.82]
%      % 72-bit word as colored blocks
%      \node[font=\scriptsize\bfseries, TealDarkCyan] at (1.8,3.5) {72-bit ECC word};
%
%      % Data bits (64): draw 4 rows of 16 → simplified as 4 wide bars
%      \foreach \row in {1,2,3,4} {
%        \pgfmathsetmacro{\yy}{3.1-(\row-1)*0.45}
%        \node[draw=TealDarkCyan, fill=TealDarkCyan!15, rounded corners=2pt,
%              minimum width=3.2cm, minimum height=0.38cm, font=\tiny]
%          at (1.2,\yy) {data bits (16 each)};
%      }
%      % Check bits (8)
%      \node[draw=AccentCyan, fill=AccentCyan!22, rounded corners=2pt,
%            minimum width=3.2cm, minimum height=0.38cm, font=\tiny\bfseries]
%        at (1.2,1.25) {8 check bits (Hamming)};
%
%      % Brace
%      \draw[TealDarkCyan, decorate,
%            decoration={brace, amplitude=4pt, raise=3pt, mirror}]
%        (2.85,1.08) -- (2.85,3.30);
%      \node[right=10pt, font=\tiny, TealDarkCyan] at (2.85,2.2) {72 bits total};
%
%      \medskip
%      \node[draw=TealDarkCyan!40, rounded corners=4pt, fill=TealDarkCyan!5,
%            minimum width=3.4cm, minimum height=0.9cm, font=\scriptsize,
%            align=center] at (1.5,0.5)
%        {Syndrome = 0: OK\\Syndrome $\ne$ 0: fix bit};
%    \end{tikzpicture}
%  \end{column}
%\end{columns}
%\end{frame}
%
%\begin{frame}{\textbf{III. Hamming Codes}}{Shortened Hamming codes}
%\small
%\defbox[Shortening]{
%Shortened Hamming codes are obtained by fixing selected coordinates to zero and deleting them.
%}
%\begin{itemize}
%  \item They are useful when exact lengths $2^r-1$ are inconvenient.
%  \item The minimum distance remains at least $3$.
%  \item Perfectness is usually lost.
%\end{itemize}
%\begin{block}{Engineering}
%Many practical memory codes are shortened or extended Hamming codes.
%\end{block}
%\end{frame}

\begin{frame}{\textbf{III. Hamming Codes}}{Hamming codes as BCH codes}
\small
\defbox[Cyclic viewpoint]{
When $n=2^r-1$, the binary Hamming code can be realized as a narrow-sense BCH code of designed distance $3$.
}
\begin{itemize}
  \item Choose a primitive $n$-th root $\alpha\in\F_{2^r}$.
  \item Let $g(X)$ be the minimal polynomial of $\alpha$ over $\F_2$.
  \item Then the cyclic code generated by $g(X)$ is a Hamming code.
\end{itemize}
\end{frame}

\begin{frame}{\textbf{III. Hamming Codes}}{The [7,4,3] cyclic form}
\small
Over $\F_2$,
\[
  X^7-1=(X+1)(X^3+X+1)(X^3+X^2+1).
\]
Taking
\[
  g(X)=X^3+X+1
\]
gives a cyclic $[7,4,3]$ Hamming code.
\begin{itemize}
  \item This is the bridge from projective checks to polynomial ideals.
  \item The next chapter makes that bridge systematic.
\end{itemize}
\end{frame}

\begin{frame}{\textbf{III. Hamming Codes}}{Classification intuition}
\small
\begin{itemize}
  \item Perfect one-error-correcting linear codes are essentially Hamming codes, apart from trivial cases.
  \item The parity-check matrix must use every possible nonzero syndrome once.
  \item Thus the construction is forced by the Hamming bound.
\end{itemize}
%\begin{block}{Mathematical taste}
Hamming codes are beautiful because optimality, decoding, and geometry are the same statement.
%\end{block}
\end{frame}

%\begin{frame}{\textbf{III. Hamming Codes}}{Exercise: distance}
%\small
%\begin{block}{Problem}
%Let $H$ contain one representative of every point of $\mathbb{P}^{r-1}(\F_q)$. Prove that $\ker H$ has distance exactly $3$.
%\end{block}
%\textbf{Hint.}
%\begin{itemize}
%  \item No one or two columns are dependent.
%  \item Any projective line contains three dependent representatives after scaling.
%  \item Translate column dependence into a codeword support.
%\end{itemize}
%\end{frame}
%
%\begin{frame}{\textbf{III. Hamming Codes}}{Exercise: perfectness}
%\small
%\begin{block}{Problem}
%Verify directly that the $q$-ary Hamming code meets the Hamming bound with equality for $t=1$.
%\end{block}
%\[
%  n=\frac{q^r-1}{q-1},\qquad k=n-r.
%\]
%\textbf{Computation.}
%\[
%  q^k(1+n(q-1))=q^{n-r}(1+q^r-1)=q^n.
%\]
%\end{frame}
%
%\begin{frame}{\textbf{III. Hamming Codes}}{Exercise: dual weights}
%\small
%\begin{block}{Problem}
%Show that every nonzero word of the binary simplex code has weight $2^{r-1}$.
%\end{block}
%\textbf{Hint.}
%For a nonzero row vector $a\in\F_2^r$, the word $aH_r$ evaluates the nonzero linear functional $x\mapsto ax^T$ on every nonzero vector $x$.
%\begin{itemize}
%  \item Exactly half of all vectors of $\F_2^r$ evaluate to $1$.
%  \item The zero vector evaluates to $0$.
%\end{itemize}
%\end{frame}

\begin{frame}{\textbf{III. Hamming Codes}}{Summary}
\small
\begin{columns}[T,onlytextwidth]
  \begin{column}{.48\textwidth}
    \defbox[Construction]{
    \begin{itemize}\small
      \item parity-check columns are projective points
      \item no repeated projective column
      \item all syndromes are used
    \end{itemize}
    }
  \end{column}
  \begin{column}{.48\textwidth}
    \defbox[Consequences]{
    \begin{itemize}\small
      \item parameters $[2^r-1,2^r-1-r,3]_2$
      \item perfect one-error correction
      \item dual simplex code
    \end{itemize}
    }
  \end{column}
\end{columns}
\end{frame}

\begin{frame}{\textbf{III. Hamming Codes}}{Transition to cyclic codes}
\small
\begin{block}{New symmetry}
Hamming codes are governed by projective geometry. Cyclic codes are governed by the cyclic shift operator.
\end{block}
\begin{itemize}
  \item Linear code: subspace of $\F_q^n$.
  \item Hamming code: subspace defined by projective parity checks.
  \item Cyclic code: subspace stable under a cyclic permutation.
\end{itemize}
\[
  \text{shift symmetry}\quad\Longleftrightarrow\quad\text{ideals in }\F_q[X]/(X^n-1).
\]
\end{frame}

